\begin{chapterbox}
    \chapter{Die wahre Macht des zweiten Schilds (WIP)}
    \label{wip-die-wahre-macht-des-zweiten-schilds}
    \az{Jahr 75}
    
    \begin{center}
        Teil E der Neuen Abenteuer

        Fortsetzung von \hyperref[wip-eine-faehrtenleserin-entscheidet-sich]{Eine Fährtenleserin entscheidet sich (WIP)}
    \end{center}

    Krahder-Prinz Undavahr stürmt mit seiner Skelettarmee die Rietburg. Als er den wunderschönen Sternenschild findet, will er den angeschlagenen Bruderschild verbrennen. Er ahnt nicht, was für eine Macht dieser Schild beinhielt. 
\end{chapterbox}



Diese Fan-Geschichte ist noch in Arbeit. Ich teile sie bereits hier, damit sie für mich nicht verloren geht, falls mein Laptop spontan verschwinden sollte, und damit sie für euch nicht völlig verloren geht, falls ich spontan verschwinden sollte. Fühlt euch in letzterem Fall vollkommen frei, hierrin zu schmökern und nach Lust und Laune weiter daran zu werkeln. Aber solange ich ein aktiver Andori bleibe, lest bitte nicht zu genau nach, was ich geplant haben könnte. LG BBB :D



\newpage
\section{Kapitel}
\az{Jahr 75}

Ken Dorr saß in der geheimen Klappe hinter dem Thron und bibberte.





\newpage
\section{Kapitel}

Dem Bruderschild entsprangen die vier Neuen \gend{Held*innen}{Helden}.







\newpage
\section{Epiloge}


\az{Jahr 75}

Nachdem die Vorhut von Chada und Thorn per Falken die Nachricht zurückgeschickt hatten, dass die Krahder das Land verlassen hatten und es relativ sicher sei, legte die Aldebaran II am Wachsamen Wald an. 

Rhega, Stinner, Kirr und der Unbekannte Krieger blieben zurück und wünschten ihnen viel Erfolg.

Die restlichen \gend{Held*innen}{Helden}, die im Norden unterwegs und am Kampf gegen die Schwarze Kogge beteiligt gewesen waren, reisten durch den Wachsamen Wald, bereit, sich dem Tross der Andori in Cavern anzuschließen.

\trennzeichen

An einem nicht weit entfernten Lagerfeuer saßen der Königswolf Lonas, Forn, Leander und Drukil. Sie aßen. Sie wollten dringend noch diese Nacht aus dem Wald hinaus, damit Drukil sich nicht verwandelte. Der Wolf führte sie in den Süden, nach Cavern, zum Tross der Andori.

\trennzeichen

Darh saß alleine auf dem niedergebrannten Hof deren Familie in Geissried. 

Da kam der Handelszwerg Garz vorbei. 

Er erzählte von seinem Plan, sich dem Tross der Andori anzuschließen, und lud Darh ein, sich ihm anzuschließen.

\trennzeichen

Aćh und Ijsdur erschienen wieder auf der andorischen Seite des Kuolema-Gebirges und machten sich auf, zum Tross der Andori zu stoßen.

\trennzeichen

Tenaya und andere Magiebegabte führten ein kompliziertes Ritual um die Zwergeneiche herum aus. Ein riesiges schwarzes Portal öffnete sich und kündigte die Ankunft Earas aus dem hohen Norden vor.

\trennzeichen

Beim Tross der Andori vor Cavern versammelten Chada, Thorn und Kram \gend{Mitglieder*innen}{Mitglieder} aller umliegenden Völker im großen Befreiungstross.

Die Neuen \gend{Held*innen}{Helden} stapften mit den von ihnen Befreiten aus dem Gebirge herunter. 

Große Augen starrten ihnen entgegen.

Fennah bereitete ihre Arme aus und proklamierte: "Wir leben noch!"



\az{Jahr 76}

In Krahd litt die Ambacu Griun als Dienerin bestimmter \gend{Krahder*innen}{Krahder}.

\trennzeichen

Doch ein ganzer Tross war unterwegs, um der Sklaverei in Krahd ein Ende zu bereiten.

Nach der Befreiung der Winterburg bereiteten sich die verschiedenen \gend{Held*innen}{Helden} auf eigene Arten auf den Angriff auf Krahd vor.

